Meeting with Yicong;

- Bootstrap procedure for the simulation will work. Being careful with which observations you randomly draw because it should not destroy the whole the time dimension.
- declustering: plot the extremal coefficient of the raw data for each station, then remove the next day's observation and plot again the extremal coefficient. the should hover around 0.8-0.9 (independence).
- after kind of each statement that I make, I need to introduce, motivate and discuss it such that the readers (supervisors) can verify the procedure that I do.
-- extremal copulas by estimating it in a pairwise manner, is what he would suggest of thinking spatial extremal dependence.
Very hard to test the parameter whether it is zero or not because you are at the boundary of the parameter space.
--> but he said that for my case, testing whether it is equal or it is different, is a simpler method.
use simulation to bootstrap the methods because you need a benchmark to assess its performance.
My supervisor who is a super good mathematician and tenured professor, did not know the F-madogram so I want to drop that one.
- he found the division between parametric and non-parametric good structure.
- he really did not care about the AI generated text and sentence structure, but he cared about the mathematical and statistical procedure that I do.
- he said bootstrap was good procedure tho.


-- the new setup:
one research question;
three parametric methods;
one canonical comparison object;
simulation comparison of the three methods;
empirical application to KNMI daily maximum wind-gust extremes.

we will use the following three models for answering the research question:
-- brown resnick
-- extremal-t max stable model
-- spatial student-(t) copula model